% Source fingerprint: d9b09bc8adcfee911a9292ed0e0848426cf83ef8ddd2694ff35e3e9a1f58a8fc
% T12: raw TeX cells preserved; no numeric highlighting.
\begin{table}[htbp]
\centering\small\renewcommand{\arraystretch}{1.18}\setlength{\tabcolsep}{4pt}
\caption[迹定理的指数、值域与右逆]{迹定理的指数、值域与右逆。表中高阶分数阶估计沿用本节的有限指数范围；边界光滑性、数据方向与指数必须同时核对。}
\label{b1:tab:visual-t12}
\begin{tabularx}{\linewidth}{@{}>{\raggedright\arraybackslash}p{3.1cm}>{\raggedright\arraybackslash}p{4.7cm}>{\raggedright\arraybackslash}X@{}}
\toprule
对象与范围 & 迹的目标 & 满射、右逆与相容性 \\
\midrule
一阶，有限 \(p\) & 本章规定区域上的边界 \(L^p\) 迹 & 端点 \(p=1\) 的 \(L^p\) 迹不自动带有下行的精确分数阶目标 \\
一阶，\(1<p<\infty\) & \(W^{1-1/p,p}\)；半空间或有界 Lipschitz 边界 & 本章构造有界线性右逆，因而是满射 \\
半空间高阶，\(u\in W^{k,p}\) & \(\gamma_j u\in W^{k-j-1/p,p}\)，\(0\le j<k\) & 本节证明迹估计与切向导数相容；没有由一阶右逆直接得到同时匹配全部数据的右逆 \\
Lipschitz 边界的法向导数 & 一阶导数的迹与本性有界法向相乘，至少得到 \(L^p\) 数据 & 本性有界法向乘子未必保持分数阶空间；折角处可失败 \\
一维区间 & 边界是两个端点 & 使用端点取值，不套用正维边界的双点差分核 \\
\bottomrule
\end{tabularx}
\end{table}
